\documentclass[11pt,a4paper]{article}
\usepackage[T1]{fontenc}
\usepackage[utf8]{inputenc}
\usepackage{lmodern,amsmath,amssymb}
\usepackage[margin=25mm]{geometry}
\usepackage[hidelinks]{hyperref}
\hypersetup{pdftitle={The mark cost is exactly $\hbar/2$, and the worst-case theorem survives as a statistical o},pdfauthor={navstokgap research working notes}}
\newcommand{\tightlist}{\setlength{\itemsep}{0pt}\setlength{\parskip}{0pt}}
\setlength{\emergencystretch}{3em}
\setcounter{secnumdepth}{0}
\begin{document}
\hypertarget{the-mark-cost-is-exactly-hbar2-and-the-worst-case-theorem-survives-as-a-statistical-one}{%
\section{\texorpdfstring{The mark cost is exactly \(\hbar/2\), and the
worst-case theorem survives as a statistical
one}{The mark cost is exactly \textbackslash hbar/2, and the worst-case theorem survives as a statistical one}}\label{the-mark-cost-is-exactly-hbar2-and-the-worst-case-theorem-survives-as-a-statistical-one}}

Reopening the question of what was wrong with the
\href{planck-gap-derivation.md}{derivation note}: almost nothing. Its
Proposition 6 claimed a positive quantum value for the mark cost
\(\kappa=\delta\Delta\), and that claim is correct. The proof was
invalid and the constant was too large by a factor of four, and the
worst-case \emph{formulation} of \(\delta\) and \(\Delta\) as support
widths is what made the quantum instance empty. Replacing support widths
by standard deviations repairs everything at once, and two theorems then
close the gap the derivation note left open.

\textbf{Theorem A.} For a mark realized by the impulsive coupling
\(\lambda\hat y\otimes\hat P_A\) of the body's position to a probe, with
the probe's pointer read afterwards, the mark's error operator and the
impulse it delivers are canonically conjugate,

\[\hat N=\frac{\hat Q_A}{\lambda}-\hat y,\qquad \hat D=-\lambda\hat P_A,
\qquad [\hat N,\hat D]=-i\hbar,\qquad\text{hence}\qquad
\delta\Delta=\Delta_\psi\hat N\cdot\Delta_\psi\hat D\ \ge\ \frac\hbar2\]

for every probe state whatever. So \(\kappa\ge\hbar/2\), exactly, with
no error--disturbance relation and no calibration argument.

\textbf{Theorem B.} In the Gaussian statistical model, for every
protocol of marks with \(\delta_j\Delta_j\ge\kappa\), of any number and
at any times, with the initial position and velocity unknown, the
optimal test's deflection obeys

\[d^2\ \le\ \frac{2F^2\tau^3}{9\,m\kappa},\qquad\text{hence deciding at
error probability }\epsilon\text{ requires}\qquad
\tau\Delta E=\frac{F^2\tau^3}{2m}\ \ge\ 9\,z_{1-\epsilon}^2\,\kappa ,\]

with \(z_{1-\epsilon}=\Phi^{-1}(1-\epsilon)\). With Theorem A this is
\(\tau\Delta E\ge\tfrac92z_{1-\epsilon}^2\hbar\), about \(12\hbar\) at
five per cent error. The proof is the derivation note's proof unchanged:
the same certificate, the same arithmetic--geometric step, the same
integration by parts to a Dirichlet energy, with the separating
hyperplane replaced by the optimal test direction. The constant \(9\)
reappears from the same Cauchy--Schwarz inequality.

Newton's premise M3 is Theorem A for his own probe: the interval of fits
\(\Lambda\) is the corpuscle's transverse position spread and \(p\) its
transverse momentum spread, so \(\kappa=\Lambda p\) is the corpuscle's
own uncertainty product, and the heuristic value \(h/2\) of the
\href{newton-mark-floor.md}{mark-floor note} differs from the theorem's
\(\hbar/2\) by \(\pi\). Section 6 sets out the single paper this
supports. Exploratory; no ledger promotion.

\hypertarget{what-was-actually-wrong-stated-exactly}{%
\subsection{1. What was actually wrong, stated
exactly}\label{what-was-actually-wrong-stated-exactly}}

The derivation note's framework defines a mark's resolution \(\delta\)
and recoil width \(\Delta\) as worst-case interval widths: the datum
confines the position to an interval of length \(\delta\) with
certainty, and the records confine the delivered impulse to an interval
of length \(\Delta\) with certainty. Three claims about that framework
are now settled.

\begin{itemize}
\tightlist
\item
  \textbf{Theorem 1 is correct and is not affected.} For every finite
  protocol with \(\delta_j\Delta_j\ge\kappa\), deciding the comparison
  needs \(F^2\tau^3>9m\kappa\).
\item
  \textbf{The worst-case framework has no quantum instance.} Certain
  position confinement forces a compactly supported conditional state,
  whose momentum distribution has full support by Paley--Wiener. So no
  quantum mark has both widths finite, \(\kappa=\infty\), and Theorem 1
  holds emptily. This is the correction that stands.
\item
  \textbf{The conclusion of the first Proposition 6 was right anyway.}
  Its assertion was that quantum kinematics gives \(\kappa\) a positive
  value of the order of \(\hbar\), and that the Galileo comparison
  therefore has a floor of order \(\hbar\) in \(\tau\Delta E\). Theorem
  A gives \(\kappa\ge\hbar/2\) and Theorem B gives the floor, so both
  are true. What failed was the route: bounding a Wasserstein-2
  calibration disturbance by half a worst-case width, when that
  disturbance is infinite for the instruments in question.
\end{itemize}

So the defect was the choice of currency, and it cost one constant and
one proof. Standard deviations are the right currency because they are
finite for quantum states, because they are what the uncertainty
relation constrains, and because the derivation note's proof never used
containment beyond a quadratic bound, as Section 3 shows.

\hypertarget{theorem-a-the-marks-error-and-its-recoil-are-conjugate}{%
\subsection{2. Theorem A: the mark's error and its recoil are
conjugate}\label{theorem-a-the-marks-error-and-its-recoil-are-conjugate}}

Take the von Neumann model of a position mark. The body has
\(\hat y,\hat p\); the probe has \(\hat Q_A,\hat P_A\) with
\([\hat Q_A,\hat P_A]=i\hbar\); the mark is the impulsive unitary

\[U=\exp\left(-\frac{i}{\hbar}\lambda\,\hat y\,\hat P_A\right),\qquad\lambda>0 .\]

\textbf{Theorem A.} In the Heisenberg picture across the mark,

\[U^\dagger\hat Q_AU=\hat Q_A+\lambda\hat y,\qquad
U^\dagger\hat pU=\hat p-\lambda\hat P_A,\qquad
U^\dagger\hat yU=\hat y,\qquad U^\dagger\hat P_AU=\hat P_A .\]

Reading the pointer \(\hat Q_A\) and estimating the body's position by
\(\hat Q_A/\lambda\), the error operator
\(\hat N=\hat Q_A/\lambda-\hat y\) and the delivered impulse
\(\hat D=\hat p_{\rm after}-\hat p_{\rm before}=-\lambda\hat P_A\)
satisfy \([\hat N,\hat D]=-i\hbar\), hence for every joint state

\[\delta\,\Delta:=\Delta\hat N\cdot\Delta\hat D\ \ge\ \frac\hbar2 .\]

\emph{Proof.} With \(A=\frac{i}{\hbar}\lambda\hat y\hat P_A\), so that
\(U^\dagger=e^{A}\), the commutators are
\([A,\hat Q_A]=\frac{i\lambda\hat y}{\hbar}[\hat P_A,\hat Q_A]=\lambda\hat y\)
and
\([A,\hat p]=\frac{i\lambda\hat P_A}{\hbar}[\hat y,\hat p]=-\lambda\hat P_A\),
both central, so the Baker--Campbell--Hausdorff series terminates and
gives the four displayed relations. The error operator is then
\(\hat N=(\hat Q_A+\lambda\hat y)/\lambda-\hat y=\hat Q_A/\lambda\),
whence

\[[\hat N,\hat D]=\left[\frac{\hat Q_A}{\lambda},-\lambda\hat P_A\right]
=-[\hat Q_A,\hat P_A]=-i\hbar,\]

and Robertson's inequality gives
\(\Delta\hat N\,\Delta\hat D\ge\frac12|\langle[\hat N,\hat D]\rangle|=\hbar/2\).
\(\square\)

Three features matter. The coupling strength \(\lambda\) cancels, so the
bound is a property of the mark and not of how hard it is made. The
result holds for every probe state, pure or mixed, Gaussian or not. And
the two quantities are the probe's own position and momentum spreads:
\(\delta=\Delta\hat Q_A/\lambda\) and \(\Delta=\lambda\Delta\hat P_A\),
so

\[\kappa=\delta\Delta=\Delta\hat Q_A\cdot\Delta\hat P_A ,\]

the mark's cost is the probe's uncertainty product, and the floor of the
Galileo comparison is set by the probe rather than by the body. The
contested error--disturbance relations play no part: Ozawa's
counterexamples concern instruments whose commutator \([\hat N,\hat D]\)
differs from \(-i\hbar\), and the present statement is confined to the
momentum-transfer class where it equals \(-i\hbar\).

\textbf{Why the records cannot recover the recoil.} The pointer record
is \(\hat R=\hat Q_A+\lambda\hat y\) and the impulse is
\(\hat D=-\lambda\hat P_A\), with \([\hat R,\hat D]=-i\hbar\lambda\).
The record and the recoil are incompatible observables, so no reading of
the pointer, however precise, determines the impulse. This is
Proposition 5 of the derivation note seen from the other side: the
far-screen construction recovers the impulse by measuring the probe's
momentum, and measuring the probe's momentum is exactly what reading its
position forbids.

\hypertarget{theorem-b-the-same-proof-in-the-gaussian-statistical-model}{%
\subsection{3. Theorem B: the same proof, in the Gaussian statistical
model}\label{theorem-b-the-same-proof-in-the-gaussian-statistical-model}}

\textbf{The model.} Marks at \(0=t_0<t_1<\dots<t_k\le\tau\). Mark \(j\)
returns \(R_j=y(t_j)+\xi_j\) with \(\xi_j\) centred Gaussian of standard
deviation \(\delta_j\), and delivers an impulse \(\iota_j\), centred
Gaussian of standard deviation \(\Delta_j\), all independent. The
trajectory is

\[y(t)=y_0+v_0t+\theta P(t)+\frac1m\sum_{j}\iota_j(t-t_j)_+,
\qquad P(t)=\frac{Ft^2}{2m},\]

with \(\theta=0\) under \(\mathrm I\) and \(\theta=1\) under
\(\mathrm F\), and \(y_0,v_0\) unknown. A protocol satisfies the
\textbf{mark trade-off} with constant \(\kappa\) when
\(\delta_j\Delta_j\ge\kappa\) for every \(j\).

\textbf{Invariance and the deflection.} Tests are required to be
invariant under the unknown \(y_0\) and \(v_0\), so the statistic is
\(u^{\mathsf T}R\) with \(u\) orthogonal to the vectors \((1,\dots,1)\)
and \((t_0,\dots,t_k)\). For Gaussian data the optimal invariant test
has error probability \(\Phi(-d/2)\) with

\[d^2=\sup\left\{\frac{(u^{\mathsf T}P)^2}{u^{\mathsf T}\Sigma u}\ :\
\sum_iu_i=0,\ \sum_iu_it_i=0\right\},\qquad P_i=P(t_i),\]

\(\Sigma\) being the covariance of \(R\). This \(u\) is the statistical
counterpart of the derivation note's separating multipliers \(\mu\), and
it obeys the same constraint \(\sum_iu_i=0\).

\textbf{Theorem B.} For every such protocol,
\(d^2\le\dfrac{2F^2\tau^3}{9\,m\kappa}\). Consequently an invariant test
with error probability at most \(\epsilon\) requires

\[F^2\tau^3\ \ge\ 18\,z_{1-\epsilon}^2\,m\kappa,\qquad
\tau\Delta E\ \ge\ 9\,z_{1-\epsilon}^2\,\kappa,\qquad
A_{\rm inertial,fall}\ \ge\ \frac{6\,z_{1-\epsilon}^2\,v\kappa}{F}.\]

\emph{Proof.} The noise decomposes exactly as in the derivation note's
feasibility system: \(\xi_j\) enters \(R_j\) alone, and \(\iota_j\)
enters \(R_i\) for \(i>j\) with coefficient \((t_i-t_j)/m\). Hence

\[u^{\mathsf T}\Sigma u=\sum_j\delta_j^2u_j^2+\sum_j\frac{\Delta_j^2}{m^2}S_j^2,
\qquad S_j=\sum_{i>j}u_i(t_i-t_j),\]

the same \(S_j\) as there. By the arithmetic--geometric mean inequality
applied termwise, exactly as in that note's Lemma 2,

\[u^{\mathsf T}\Sigma u\ \ge\ \sum_j\frac{2\delta_j\Delta_j}{m}|u_j||S_j|
\ \ge\ \frac{2\kappa}{m}\sum_j|u_j||S_j| .\]

Since \(\sum_iu_i=0\), Lemma 3 of that note applies verbatim with
\(\mu\) replaced by \(u\): with
\(T(w)=\int_w^\tau\sum_{i:t_i>s}u_i\,ds\) one has \(T(t_j)=S_j\),

\[u^{\mathsf T}P=\frac Fm\int_0^\tau T,\qquad
\sum_ju_jT(t_j)=-E,\qquad E=\int_0^\tau T'^2 ,\]

so \(\sum_j|u_j||S_j|\ge E\) and
\(u^{\mathsf T}\Sigma u\ge2\kappa E/m\). Cauchy--Schwarz with
\(T(\tau)=0\) gives \(\int_0^\tau T\le\frac23\tau^{3/2}\sqrt E\), so

\[\frac{(u^{\mathsf T}P)^2}{u^{\mathsf T}\Sigma u}
\le\frac{(F/m)^2\frac49\tau^3E}{2\kappa E/m}=\frac{2F^2\tau^3}{9m\kappa},\]

independently of \(u\), of the number of marks and of their times. The
error probability \(\Phi(-d/2)\le\epsilon\) requires
\(d\ge2z_{1-\epsilon}\), and
\(F^2\tau^3\ge\frac92m\kappa d^2\ge18z_{1-\epsilon}^2m\kappa\).
\(\square\)

The worst-case theorem and the statistical theorem therefore have one
proof. The certificate changes meaning, from a hyperplane separating a
feasible set to the direction of the optimal test, and every step after
it is identical. This is the sense in which the derivation note's
content survived its formulation: the argument never needed interval
containment, only a quadratic bound in the same two variables.

\hypertarget{why-gaussian-quantum-experiments-are-covered-exactly}{%
\subsection{4. Why Gaussian quantum experiments are covered
exactly}\label{why-gaussian-quantum-experiments-are-covered-exactly}}

Within its class the statistical model of Section 3 reproduces the
quantum one exactly. For Gaussian probe states, linear dynamics and
pointer readouts, the Wigner function is a genuine probability density
and evolves by the classical linear equations, so every measured
distribution in the quantum experiment equals the corresponding
distribution in the classical Gaussian model whose noise covariance is
the probe's Wigner covariance. Theorem A identifies that covariance's
uncertainty product as \(\delta\Delta\ge\hbar/2\). Hence for every
protocol of Gaussian momentum-transfer marks on a body under the two
hypotheses,

\[\tau\Delta E\ \ge\ \frac92\,z_{1-\epsilon}^2\,\hbar .\]

At \(\epsilon=0.05\), where \(z=1.645\), this is
\(\tau\Delta E\ge12.2\,\hbar\); at \(\epsilon=0.32\), where \(z=1\), it
is \(4.5\,\hbar\). The insertion statement follows as before: marks
inserted into a window of duration \(\tau'\), used by themselves, record
the force only if \(F^2\tau'^3\ge18z^2m\kappa\), so Democritus insertion
stops at the mesh

\[\tau_*=\left(\frac{18\,z_{1-\epsilon}^2\,m\kappa}{F^2}\right)^{1/3}
=\left(\frac{9\,z_{1-\epsilon}^2\,m\hbar}{F^2}\right)^{1/3},\]

and finer marks record free motion at the stated confidence.

This coexists with the counterexample of the
\href{planck-gap-probabilistic.md}{probabilistic note}, which
distinguishes the hypotheses at arbitrarily small \(\tau\Delta E\) using
a two-packet preparation of separation \(\pi\hbar/(F\tau)\). That
protocol is non-Gaussian and makes no marks: it prepares once, waits,
and measures once. The two results divide the ground cleanly.
\textbf{Marking the trajectory costs \(\hbar/2\) per mark and gives a
floor of order \(\hbar\); declining to mark it costs an apparatus of
size \(\pi\hbar/(F\tau)\) and gives the floor \(\hbar^2/(4A)\).}
Newton's refinement is the first case, which is why the floor is the
relevant statement for the insertion question.

\hypertarget{newtons-premise-is-robertsons-inequality-for-his-corpuscle}{%
\subsection{5. Newton's premise is Robertson's inequality for his
corpuscle}\label{newtons-premise-is-robertsons-inequality-for-his-corpuscle}}

The \href{newton-mark-floor.md}{mark-floor note} posited M3, that a mark
fixing the position within \(\delta\) leaves the delivered impulse
undetermined within \(\Lambda p/\delta\), and observed that Newton's
optics supplies \(\Lambda\) and \(p\) while his determinism denies the
premise. Theorem A says what M3 is. The corpuscle is the probe;
\(\hat Q_A\) is its transverse position and \(\hat P_A\) its transverse
momentum; the interval of fits \(\Lambda\) is the scale of the first and
the corpuscular impulse \(p\) the scale of the second; and the premise
asserts precisely that their product cannot be reduced. So

\[\kappa=\Lambda p\quad\longleftrightarrow\quad
\kappa=\Delta\hat Q_A\cdot\Delta\hat P_A\ \ge\ \frac\hbar2 ,\]

and M3 is Robertson's inequality for the corpuscle, stated in the two
quantities Newton had. Newton measured the first as \(1/89000\) inch,
had no access to the second, held both to be determinate properties of
the corpuscle, and thereby denied exactly the inequality. The modern
value \(\hbar/2\) against the identification \(\Lambda p=h/2\) differ by
\(\pi\).

The historical claim this licenses is narrow and checkable. Newton's
optics contains both factors of the product that bounds the recorded
Galileo comparison; his mathematics contains the comparison and takes it
to zero; and the single proposition that would join them is a claim of
indeterminacy about the fits, which he considered and rejected in favour
of a determinate periodic disposition. The gap between the
\emph{Principia}'s vanishing sagitta and the \emph{Opticks}' finite
interval of fits follows from that single commitment, and the separation
of the two books is incidental to it.

\hypertarget{the-single-paper-valid-for-both-audiences}{%
\subsection{6. The single paper, valid for both
audiences}\label{the-single-paper-valid-for-both-audiences}}

The user's requirement is one paper that a foundations-of-physics
referee and a history-and-philosophy-of-science referee both accept.
That is attainable, because the historical analysis now does
argumentative work rather than decorating a theorem. The structure is:

\begin{enumerate}
\def\labelenumi{\arabic{enumi}.}
\tightlist
\item
  \textbf{The comparison and its limit.} Newton's Lemmas X and XI and
  the projectile Scholium, with the sagitta, the inertial--parabola area
  and the identity \((3F/v)A=\tau\Delta E\). Berkeley's objection and
  its resolution inside the geometry, with Guicciardini on what the
  limit arguments were for.
\item
  \textbf{The two theorems.} Theorem A, that a momentum-transfer mark's
  error and recoil are conjugate, and Theorem B, that every protocol of
  such marks needs \(\tau\Delta E\ge9z^2\kappa\). Proof by certificate,
  arithmetic--geometric mean and Dirichlet energy. This is the
  foundations content and it stands alone.
\item
  \textbf{The two factors in the \emph{Opticks}.} The interval of fits
  as a measured length, the corpuscular impulse as a posited one, and
  Shapiro's account of what the theory of fits was and how determinate
  Newton meant it. This is the history content and it stands alone.
\item
  \textbf{The junction.} M3 is Robertson for the corpuscle. The
  counterfactual is replaced by an exact statement: the premise that the
  theorem needs is one Newton formulated the negation of, so the
  distance between his system and a positive \(h\) is a single
  proposition, nameable in his own vocabulary.
\item
  \textbf{What the theorem does not give.} The aperture counterexample,
  the standard-quantum-limit literature and Ozawa's class, and the
  honest prior-art label: the combination \(F^2\tau^3\gtrsim m\hbar\) is
  the standard quantum limit, and the contribution is universality over
  protocols, the exact mark cost, and the identification of the premise.
\end{enumerate}

The venues that take this shape are Foundations of Physics and Studies
in History and Philosophy of Modern Physics, both of which publish
papers whose historical and technical halves are load-bearing for each
other. Two obligations remain before submission: reading Shapiro,
Guicciardini and Sabra rather than citing them, and replacing the
Project Gutenberg \emph{Opticks} and the Wikisource \emph{Principia}
with the fourth edition and the Cohen--Whitman translation.

\hypertarget{consequence-for-state}{%
\subsection{7. Consequence for STATE}\label{consequence-for-state}}

The derivation note's Proposition 6 is reinstated in substance with a
correct proof and the constant \(\hbar/2\) in place of \(2\hbar\), and
its worst-case formulation is replaced by standard deviations. The
Planck gap for marked trajectories is
\(\tau\Delta E\ge\frac92z_{1-\epsilon}^2\hbar\), protocol-universal,
with the insertion mesh \(\tau_*=(9z^2m\hbar/F^2)^{1/3}\). The remaining
programme is item 6's two obligations plus: extending Theorem A beyond
the momentum-transfer class, where Ozawa's instruments live; the general
force law, which Theorem B's proof already reduces to a norm of \(P''\);
and closing the constants between the worst-case, statistical and
single-shot forms.
\end{document}
